\section{Bốn đường dựng ground truth và cái bị ép}
\label{sec:ch16-bon-duong-dung-san}

\index{ground truth!bốn đường, gom lại|(}%
Câu hỏi mới mà mục trước để lại là câu hỏi về cái giá, và với đường thứ nhất thì
không hỏi được nếu chưa gom đủ các dạng của nó. Sách đã gặp bốn dạng, ở bốn chỗ
cách nhau bảy chương, và mục này đặt cả bốn cạnh nhau lần đầu.

\index{ground truth!đường thứ nhất, thiết kế nhiệm vụ}%
Đường thứ nhất ép bằng thiết kế nhiệm vụ.
Chương~\ref{ch:chi-so-khong-do-do-trung-thuc} đọc nó: chọn những bài toán mà lời
giải đúng buộc phải đi qua một bước nhất định, rồi đọc bước ấy ra từ chính đáp
án mà mô hình đưa. Người dựng là người muốn kiểm định chỉ số, và họ trả giá bằng
việc chỉ được chấm điểm trên những bài toán có cấu trúc ấy.

\index{ground truth!đường thứ hai, thiết kế dữ liệu}%
Đường thứ hai ép bằng thiết kế dữ liệu. Bài của
mục~\ref{sec:ch14-sua-bang-nhan-qua} tự viết các phương trình sinh dữ liệu, nên
quan hệ giữa các biến là quan hệ mà chính người thí nghiệm đặt vào, và một tập
giá trị tham chiếu tồn tại vì có người viết ra nó. Người dựng ở đây cũng là
người muốn đo, và cái bị ép không phải nhiệm vụ mà là phân phối.

\index{ground truth!đường thứ ba, thiết kế mô hình}%
Đường thứ ba ép bằng thiết kế chính mô hình, và
mục~\ref{sec:ch14-ground-truth-ke-tan-cong} đã đọc kỹ nó vì người dựng ở đó
không giống hai đường trên. Mô hình thiên lệch trong thí nghiệm
\tn{phép dựng giàn}{scaffolding} được viết ra để chỉ dự đoán bằng thuộc tính
được bảo vệ, nên trên mọi điểm dữ liệu thật thì câu trả lời trung thực đã biết
trước. Người dựng là bên tấn công, và họ dựng không phải vì muốn kiểm định gì:
họ buộc phải biết mô hình của mình dựa vào cái gì thì mới giấu được cái ấy.

\index{ground truth!đường thứ tư, nhãn có sẵn}%
Đường thứ tư thì mục~\ref{sec:ch15-dung-cu-duoc-kiem} vừa thêm vào, và nó khác
ba đường trên ở chỗ không ai dựng gì cả. Trong một bộ dữ liệu có gán nhãn khái
niệm, lượng thông tin mà mỗi khái niệm thật mang về nhãn nhiệm vụ là một đại
lượng tính thẳng từ bảng nhãn, không phụ thuộc vào mô hình nào. Người gán nhãn
đã làm việc ấy từ trước, cho một mục đích khác, và cái mà
định nghĩa~\ref{def:ch15-hai-diem-so} cần đã nằm sẵn ở đó.
Hình~\ref{fig:ch16-bon-duong-dung-san} đặt cả bốn cạnh nhau.

\begin{figure}[htbp]
  \centering
  \begin{tikzpicture}[
  every node/.style={font=\footnotesize, align=center},
  known/.style={draw=black!85, line width=1pt, fill=black!12,
    minimum width=30mm, minimum height=9mm},
  unknown/.style={draw=black!55, dotted, thick, fill=white,
    minimum width=34mm, minimum height=9mm},
  route/.style={draw=black!65, rounded corners=2pt, fill=black!5,
    minimum width=26mm, minimum height=15mm, text width=24mm,
    font=\scriptsize},
  flow/.style={-{Stealth[length=2mm]}, thick, black!70},
  thinflow/.style={-{Stealth[length=1.6mm]}, black!55},
  lbl/.style={font=\scriptsize, align=center, black!55}
]
  \node[known]   (gt)  at (1.4,4.2) {lời giải thích};
  \node[unknown] (dep) at (8.6,4.2) {cái mà hành vi\\phụ thuộc vào};

  \draw[flow] (gt.east) -- (dep.west);
  \node[lbl] at (5.0,4.85) {định nghĩa~\ref{def:ch01-do-trung-thuc} so hai bên,
    dưới một tập can thiệp đã nêu};

  \node[route] (r1) at (0,0)   {thiết kế \textbf{nhiệm vụ}\\[1mm] câu trả lời đúng buộc phải đi qua một bước};
  \node[route] (r2) at (3.1,0) {thiết kế \textbf{dữ liệu}\\[1mm] quan hệ giữa các biến do phương trình sinh đặt ra};
  \node[route] (r3) at (6.2,0) {thiết kế \textbf{mô hình}\\[1mm] mô hình được viết để chỉ dùng một đặc trưng};
  \node[route] (r4) at (9.3,0) {\textbf{gán nhãn} sẵn\\[1mm] không ai dựng gì; nhãn khái niệm đã nằm sẵn trong dữ liệu};

  \draw[thinflow] (r1.north) -- (0,1.5);
  \draw[thinflow] (r2.north) -- (3.1,1.5);
  \draw[thinflow] (r3.north) -- (6.2,1.5);
  \draw[thinflow] (r4.north) -- (9.3,1.5);
  \draw[black!55] (0,1.5) -- (9.3,1.5);
  \draw[flow] (8.6,1.5) -- (dep.south);
  \node[lbl] at (6.35,2.35) {bốn cách cấp vế phải};

  \node[lbl] at (0,-1.35)   {chương~\ref{ch:chi-so-khong-do-do-trung-thuc}};
  \node[lbl] at (3.1,-1.35) {mục~\ref{sec:ch14-sua-bang-nhan-qua}};
  \node[lbl] at (6.2,-1.35) {mục~\ref{sec:ch14-ground-truth-ke-tan-cong}};
  \node[lbl] at (9.3,-1.35) {chương~\ref{ch:nut-that-khai-niem}};
\end{tikzpicture}

  \caption{Bốn đường không khác nhau ở chỗ chúng cấp cái gì, vì cả bốn đều cấp
  đúng một thứ, là vế phải mà định nghĩa~\ref{def:ch01-do-trung-thuc} không có.
  Chúng khác nhau ở chỗ cái gì phải bị ép để vế ấy biết được, và ba đường đầu ép
  một thứ mà đường thứ tư không phải ép, vì thứ ấy đã có sẵn. Ô bên phải vẽ nét
  chấm theo quy ước của sách cho thứ vắng mặt: với một mô hình bắt được ngoài
  đời thì không ai biết nội dung của nó, và đó là chỗ mắc kẹt mà cả bốn đường
  tồn tại để vòng qua.}
  \label{fig:ch16-bon-duong-dung-san}
\end{figure}

\index{ground truth!chỗ chung của bốn đường}%
Đặt cạnh nhau thì thấy hai điều mà đọc rời từng đường không thấy. Điều thứ nhất
là cả bốn đường đều cấp đúng một thứ: vế phải của định
nghĩa~\ref{def:ch01-do-trung-thuc}, thứ mà với một mô hình có sẵn thì không ai
biết. Chúng không phải bốn cách đo khác nhau mà là bốn cách trả lời cùng một câu
hỏi thiếu, nên so chúng với nhau là việc làm được.

Điều thứ hai là ba đường đầu trả cùng một kiểu giá và đường thứ tư trả kiểu
khác. Ở ba đường đầu, cái làm cho vế phải biết được cũng chính là cái làm cho
kết luận hẹp lại: nhiệm vụ được chọn, phân phối được viết, mô hình được dựng
riêng, và kết luận rút ra là kết luận về nhiệm vụ ấy, phân phối ấy, mô hình ấy.
Mục~\ref{sec:ch14-ground-truth-ke-tan-cong} đã phát biểu điều đó cho đường thứ
ba trong một câu đáng lặp lại ở đây, vì bây giờ nó áp cho ba đường: một phép dựng
cho ta ground truth và lấy đi tính đại diện, đổi cái này lấy cái kia.

\index{ground truth!cái giá của đường thứ tư}%
Đường thứ tư không đổi như vậy, vì nó không dựng gì để mà đánh đổi. Cái nó đòi
là một điều kiện có trước: bộ dữ liệu phải đã được gán nhãn khái niệm, và tập
nhãn ấy phải đủ dự báo nhiệm vụ. Hai vế ấy đều đắt theo nghĩa thông thường, và
mục~\ref{sec:ch15-nguyen-nhan} đã cho thấy vế thứ hai hỏng thì hỏng ra sao:
chính người đề xuất khung viết rằng khi không kiếm thêm được nhãn khái niệm thì
phải kết luận nhiệm vụ ấy không dùng lối tiếp cận theo khái niệm được. Nhưng cái
giá ấy là giá phải trả trước khi bắt đầu, không phải chỗ hẹp trong kết luận rút
ra sau. Trong bốn đường thì đây là đường duy nhất mà kết luận nói về một bộ dữ
liệu thật, và cũng là đường duy nhất mà phần lớn các nhiệm vụ không đi được.

\index{ground truth!bốn đường, gom lại|)}%
Bốn đường ấy là bốn dạng của một lối thoát, nên cái giá của lối thoát ấy phải
đọc trên dạng nào đo được nó. Ba dạng sau chưa dạng nào có bài báo tự đếm xem
mình với tới đâu. Dạng thứ nhất thì có, vì bài của
chương~\ref{ch:chi-so-khong-do-do-trung-thuc} in ra một bảng phân bố nhãn cho
phép đếm, và mục sau đọc bảng ấy.
